\section{Casi d'Uso}

\subsection{Attori}
\begin{itemize}
	\item \textbf{Utente}: Attore primario. Rappresenta l'utilizzatore finale che interagisce con l'interfaccia dell'editor.
	\item \textbf{LLM (Large Language Model)}: Attore secondario (sistema esterno). Fornisce i servizi di analisi e generazione testo.
\end{itemize}

% --------------------------------------------------------------------------
% 1. Macroarea 1: Gestione Editor e Markdown
% --------------------------------------------------------------------------
\newpage
\subsection{Macroarea 1 — Gestione Editor e Markdown}
\textit{Funzionalità core per la manipolazione e visualizzazione del testo.}

% UC1 - Creazione di una nuova nota vuota
\subsubsection{UC 1: Creazione di una nuova nota vuota}

% use case diagram
\begin{figure}[H]
    \centering
    \includegraphics[width=0.6\textwidth]{creazione_nota.png}
    \caption{UC1 – Creazione di una nuova nota}
\end{figure}

\begin{itemize}
    \item \textbf{Attore Principale}: Utente
    \item \textbf{Scenario Principale}: 
    \begin{enumerate}
        \item L'utente attiva il comando "Nuova Nota";
        \item Il sistema crea una nuova nota vuota.
    \end{enumerate}
    \item \textbf{Precondizioni}: Il sistema è avviato e l'interfaccia è correttamente caricata.
    \item \textbf{Postcondizioni}: Viene creata una nota vuota.
    \item \textbf{Trigger}: L'utente desidera iniziare una nuova sessione di scrittura.
\end{itemize}

\subsubsection{UC 2: Eliminazione di una nota}

% use case diagram
\begin{figure}[H]
    \centering
    \includegraphics[width=0.6\textwidth]{eliminazione_nota.png}
    \caption{UC2 – Eliminazione di una nota}
\end{figure}

\begin{itemize}
    \item \textbf{Attore Principale}: Utente
    \item \textbf{Scenario Principale}: L'utente elimina la nota selezionata.
    \item \textbf{Precondizioni}: Il sistema è avviato ed è presente una nota.
    \item \textbf{Postcondizioni}: La nota è stata elimintata.
    \item \textbf{Trigger}: L'utente decide di rimuovere una nota non più necessaria.
\end{itemize}

% UC2.1 - Conferma eliminazione
\paragraph{UC 2.1: Conferma eliminazione}\mbox{}\

% use case diagram
\begin{figure}[H]
    \centering
    \includegraphics[width=0.6\textwidth]{conferma_eliminazione.png}
    \caption{UC2.1 – Conferma eliminazione}
\end{figure}

\begin{itemize}
    \item \textbf{Attore Principale}: Utente
    \item \textbf{Scenario Principale}: l'utente conferma l'eliminazione della nota.
    \item \textbf{Precondizioni}: L'utente ha selezionato una nota da eliminare.
    \item \textbf{Postcondizioni}: La nota è stata eliminata.
    \item \textbf{Trigger}: L'utente da conferma all'eliminazione di una nota non più necessaria.
\end{itemize}

% UC2.2 - Annullamento eliminazione
\paragraph{UC 2.2: Annullamento eliminazione}\mbox{}\

% use case diagram
\begin{figure}[H]
    \centering
    \includegraphics[width=0.6\textwidth]{annullamento_eliminazione.png}
    \caption{UC2.2 – Annullamento eliminazione}
\end{figure}

\begin{itemize}
    \item \textbf{Attore Principale}: Utente
    \item \textbf{Scenario Principale}: l'utente annulla l'eliminazione della nota.
    \item \textbf{Precondizioni}: L'utente ha selezionato una nota da eliminare.
    \item \textbf{Postcondizioni}: La nota non è stata eliminata.
    \item \textbf{Trigger}: L'utente decide di annullare l'eliminazione di una nota.
\end{itemize}

% UC3 - Scrittura di testo nell'editor
\subsubsection{UC 3: Scrittura testo nell'editor}

% use case diagram
\begin{figure}[H]
    \centering
    \includegraphics[width=0.6\textwidth]{scrittura_nota.png}
    \caption{UC3 – Scrittura di testo nell'editor}
\end{figure}

\begin{itemize}
    \item \textbf{Attore Principale}: Utente
    \item \textbf{Scenario Principale}: 
    \begin{enumerate}
        \item L'utente inserisce testo nell'area di input tramite tastiera.
    \end{enumerate}
    \item \textbf{Precondizioni}: L'editor è attivo e pronto a ricevere input.
    \item \textbf{Postcondizioni}: Il testo inserito è visibile nell'area di input.
    \item \textbf{Trigger}: L'utente ha bisogno di inserire o modificare i contenuti testuali del documento.
\end{itemize}

% UC4 - Copia testo
\subsubsection{UC 4: Copia testo}

% use case diagram
\begin{figure}[H]
    \centering
    \includegraphics[width=0.6\textwidth]{copia_testo.png}
    \caption{UC4 – Copia testo}
\end{figure}

\begin{itemize}
    \item \textbf{Attore Principale}: Utente
    \item \textbf{Scenario Principale}: 
    \item \begin{enumerate}
        \item L'utente seleziona una porzione di testo.
        \item L'utente attiva il comando di copia.
        \item Il sistema copia il testo selezionato negli appunti.
    \end{enumerate}
    \item \textbf{Precondizioni}: L'editor è attivo ed una porzione di testo è selezionata.
    \item \textbf{Postcondizioni}: Il testo è stato copiato.
    \item \textbf{Trigger}: L'utente vuole copiare una porzione di testo.
\end{itemize}

% UC5 - Taglia testo
\subsubsection{UC 5: Taglia testo}

% use case diagram
\begin{figure}[H]
    \centering
    \includegraphics[width=0.6\textwidth]{taglia_testo.png}
    \caption{UC5 – Taglia testo}
\end{figure}

\begin{itemize}
    \item \textbf{Attore Principale}: Utente
    \item \textbf{Scenario Principale}: 
    \begin{enumerate}
        \item L'utente seleziona il testo da tagliare.
        \item L'utente attiva il comando di taglio.
        \item Il sistema taglia il testo selezionato e lo inserisce negli appunti.
    \end{enumerate}
    \item \textbf{Precondizioni}: L'editor è attivo ed una porzione di testo è selezionata.
    \item \textbf{Postcondizioni}: Il testo è stato tagliato.
    \item \textbf{Trigger}: L'utente vuole tagliare una porzione di testo.
\end{itemize}

% UC6 - Incolla testo
\subsubsection{UC 6: Incolla testo}

% use case diagram
\begin{figure}[H]
    \centering
    \includegraphics[width=0.6\textwidth]{incolla_testo.png}
    \caption{UC6 – Incolla testo}
\end{figure}

\begin{itemize}
    \item \textbf{Attore Principale}: Utente
    \item \textbf{Scenario Principale}: 
    \begin{enumerate}
        \item L'utente seleziona dove incollare il testo.
        \item L'utente attiva il comando di incolla.
        \item Il sistema incolla il testo nella posizione desiderata.
    \end{enumerate}
    \item \textbf{Precondizioni}: L'editor è attivo ed una porzione di testo è stata copiata o tagliata.
    \item \textbf{Postcondizioni}: Il testo è stato incollato.
    \item \textbf{Trigger}: L'utente vuole incollare una porzione di testo.
\end{itemize}

% UC7 - Attivazione grassetto
\subsubsection{UC 7: Attivazione grassetto}

% use case diagram
\begin{figure}[H]
    \centering
    \includegraphics[width=0.6\textwidth]{attivazione_grassetto.png}
    \caption{UC7 – Attivazione grassetto}
\end{figure}

\begin{itemize}
    \item \textbf{Attore Principale}: Utente
    \item \textbf{Scenario Principale}: L'utente seleziona lo stile grassetto per scrivere il testo.
    \item \textbf{Precondizioni}: L'editor è attivo.
    \item \textbf{Postcondizioni}: Il testo inserito sarà in grassetto.
    \item \textbf{Trigger}: L'utente vuole applicare lo stile grassetto al testo successivo.
\end{itemize}

% UC7.1 - Grassetto su testo selezionato
\paragraph{UC 7.1: Grassetto su testo selezionato}\mbox{}\

% use case diagram
\begin{figure}[H]
    \centering
    \includegraphics[width=0.6\textwidth]{grassetto_testo_selezionato.png}
    \caption{UC7.1 – Grassetto su testo selezionato}
\end{figure}

\begin{itemize}
    \item \textbf{Attore Principale}: Utente
    \item \textbf{Scenario Principale}: L'utente seleziona del testo e applica lo stile grassetto.
    \item \textbf{Precondizioni}: L'editor è attivo e una porzione di testo è selezionata.
    \item \textbf{Postcondizioni}: Il testo selezionato sarà in grassetto.
    \item \textbf{Trigger}: L'utente vuole applicare lo stile grassetto al testo selezionato.
\end{itemize}

% UC8 -  Disattivazione grassetto
\subsubsection{UC 8: Disattivazione grassetto}

% use case diagram
\begin{figure}[H]
    \centering
    \includegraphics[width=0.6\textwidth]{disattivazione_grassetto.png}
    \caption{UC8 – Disattivazione grassetto}
\end{figure}

\begin{itemize}
    \item \textbf{Attore Principale}: Utente
    \item \textbf{Scenario Principale}: L'utente disattiva lo stile grassetto nella scrittura del testo.
    \item \textbf{Precondizioni}: L'editor è attivo.
    \item \textbf{Postcondizioni}: Il testo inserito non sarà in grassetto.
    \item \textbf{Trigger}: L'utente vuole disattivare lo stile grassetto.
\end{itemize}

% UC8.1 - Disattivazione grassetto su testo selezionato
\paragraph{UC 8.1: Disattivazione grassetto su testo selezionato}\mbox{}\

% use case diagram
\begin{figure}[H]
    \centering
    \includegraphics[width=0.6\textwidth]{disattivazione_grassetto_testo_selezionato.png}
    \caption{UC8.1 – Disattivazione grassetto su testo selezionato}
\end{figure}

\begin{itemize}
    \item \textbf{Attore Principale}: Utente
    \item \textbf{Scenario Principale}: L'utente seleziona del testo e applica lo stile grassetto.
    \item \textbf{Precondizioni}: L'editor è attivo e una porzione di testo è selezionata.
    \item \textbf{Postcondizioni}: Il testo selezionato sarà in grassetto.
    \item \textbf{Trigger}: L'utente vuole applicare lo stile grassetto al testo selezionato.
\end{itemize}

% UC9 -  Attivazione corsivo
\subsubsection{UC 9: Attivazione corsivo}

% use case diagram
\begin{figure}[H]
    \centering
    \includegraphics[width=0.6\textwidth]{attivazione_corsivo.png}
    \caption{UC9 – Attivazione corsivo}
\end{figure}

\begin{itemize}
    \item \textbf{Attore Principale}: Utente
    \item \textbf{Scenario Principale}: L'utente attiva lo stile corsivo nella scrittura del testo.
    \item \textbf{Precondizioni}: L'editor è attivo.
    \item \textbf{Postcondizioni}: Il testo inserito sarà in corsivo.
    \item \textbf{Trigger}: L'utente vuole applicare lo stile corsivo al testo successivo.
\end{itemize}

% UC9.1 - Attivazione corsivo su testo selezionato
\paragraph{UC 9.1: Attivazione corsivo su testo selezionato}\mbox{}\

% use case diagram
\begin{figure}[H]
    \centering
    \includegraphics[width=0.6\textwidth]{attivazione_corsivo_testo_selezionato.png}
    \caption{UC9.1 – Attivazione corsivo su testo selezionato}
\end{figure}

\begin{itemize}
    \item \textbf{Attore Principale}: Utente
    \item \textbf{Scenario Principale}: L'utente seleziona del testo e applica lo stile corsivo.
    \item \textbf{Precondizioni}: L'editor è attivo e una porzione di testo è selezionata.
    \item \textbf{Postcondizioni}: Il testo selezionato sarà in corsivo.
    \item \textbf{Trigger}: L'utente vuole applicare lo stile corsivo al testo selezionato.
\end{itemize}

% UC10 -  Disattivazione corsivo
\subsubsection{UC 10: Disattivazione corsivo}

% use case diagram
\begin{figure}[H]
    \centering
    \includegraphics[width=0.6\textwidth]{disattivazione_corsivo.png}
    \caption{UC10 – Disattivazione corsivo}
\end{figure}

\begin{itemize}
    \item \textbf{Attore Principale}: Utente
    \item \textbf{Scenario Principale}: L'utente disattiva lo stile corsivo nella scrittura del testo.
    \item \textbf{Precondizioni}: L'editor è attivo.
    \item \textbf{Postcondizioni}: Il testo inserito non sarà in corsivo.
    \item \textbf{Trigger}: L'utente vuole disattivare lo stile corsivo.
\end{itemize}

% UC10.1 - Disattivazione corsivo su testo selezionato
\paragraph{UC 10.1: Disattivazione corsivo su testo selezionato}\mbox{}\

% use case diagram
\begin{figure}[H]
    \centering
    \includegraphics[width=0.6\textwidth]{disattivazione_corsivo_testo_selezionato.png}
    \caption{UC10.1 – Disattivazione corsivo su testo selezionato}
\end{figure}

\begin{itemize}
    \item \textbf{Attore Principale}: Utente
    \item \textbf{Scenario Principale}: L'utente seleziona del testo e disattiva lo stile corsivo.
    \item \textbf{Precondizioni}: L'editor è attivo e una porzione di testo è selezionata.
    \item \textbf{Postcondizioni}: Il testo selezionato non sarà in corsivo.
    \item \textbf{Trigger}: L'utente vuole disattivare lo stile corsivo.
\end{itemize}

% UC5 - Inserimento elementi Markdown
\subsubsection{UC 5: Inserimento elementi Markdown}

% use case diagram
\begin{figure}[H]
    \centering
    \includegraphics[width=0.6\textwidth]{inserimento_markdown.png}
    \caption{UC5 – Inserimento elementi Markdown}
\end{figure}

\begin{itemize}
    \item \textbf{Attore Principale}: Utente
    \item \textbf{Scenario Principale}: 
    \begin{enumerate}
        \item L'utente seleziona la funzione di markdown.
        \item L'utente applica la formattazione desiderata.
    \end{enumerate}
    \item \textbf{Precondizioni}: Viene selezionata una funzione di markdown.
    \item \textbf{Postcondizioni}: Il testo viene formattato con quella funzione.
    \item \textbf{Trigger}: L'utente vuole strutturare il documento aggiungendo blocchi specifici del Markdown.
\end{itemize}

% UC5.1 - Inserimento titoli/intestazioni
\paragraph{UC 5.1: Inserimento titoli/intestazioni}\mbox{}\

% use case diagram
\begin{figure}[H]
    \centering
    \includegraphics[width=0.6\textwidth]{inserimento_markdown_intestazioni.png}
    \caption{UC5.1 – Inserimento titoli/intestazioni}
\end{figure}

\begin{itemize}
    \item \textbf{Attore Principale}: Utente
    \item \textbf{Scenario Principale}: 
    \begin{enumerate}
        \item L'utente seleziona la funzione di intestazione.
        \item L'utente applica l'intestazione desiderata secondo la sintassi di Markdown.
    \end{enumerate}
    \item \textbf{Precondizioni}: Viene selezionata la formattazione per le intestazioni.
    \item \textbf{Postcondizioni}: Il testo presenta le intestazioni.
    \item \textbf{Trigger}: L'utente vuole creare una nuova sezione o gerarchia nel testo inserendo un'intestazione.
\end{itemize}

% UC5.2 - Inserimento link ipertestuale
\paragraph{UC 5.2: Inserimento link ipertestuale}\mbox{}\

% use case diagram
\begin{figure}[H]
    \centering
    \includegraphics[width=0.6\textwidth]{inserimento_markdown_link.png}
    \caption{UC5.2 – Inserimento link ipertestuale}
\end{figure}

\begin{itemize}
    \item \textbf{Attore Principale}: Utente
    \item \textbf{Scenario Principale}: 
    \begin{enumerate}
        \item L'utente seleziona la funzione di link.
        \item L'utente inserisce l'URL e il testo del link.
    \end{enumerate}
    \item \textbf{Precondizioni}: Viene selezionata la formattazione per i link.
    \item \textbf{Postcondizioni}: Il testo presenta i link.
    \item \textbf{Trigger}: L'utente vuole inserire un collegamento verso una pagina web o una risorsa esterna.
\end{itemize}

% UC5.3 - Inserimento elenchi puntati e numerati
\paragraph{UC 5.3: Inserimento elenchi puntati e numerati}\mbox{}\

% use case diagram
\begin{figure}[H]
    \centering
    \includegraphics[width=0.6\textwidth]{inserimento_markdown_elenchi.png}
    \caption{UC5.3 – Inserimento elenchi puntati e numerati}
\end{figure}

\begin{itemize}
    \item \textbf{Attore Principale}: Utente
    \item \textbf{Scenario Principale}: 
    \begin{enumerate}
        \item L'utente seleziona la funzione per inserire elenchi.
        \item L'utente inserisce un elenco.
    \end{enumerate}
    \item \textbf{Precondizioni}: Viene selezionata la formattazione per gli elenchi.
    \item \textbf{Postcondizioni}: Il testo presenta gli elenchi.
    \item \textbf{Trigger}: L'utente vuole raggruppare una serie di informazioni sotto forma di lista.
\end{itemize}

% UC6 - Gestione delle modifiche
\subsubsection{UC 6: Gestione delle modifiche}

% use case diagram
\begin{figure}[H]
    \centering
    \includegraphics[width=0.6\textwidth]{gestione_modifiche.png}
    \caption{UC6 – Gestione delle modifiche}
\end{figure}

\begin{itemize}
    \item \textbf{Attore Principale}: Utente
    \item \textbf{Scenario Principale}: 
    \begin{enumerate}
        \item L'utente seleziona la funzione di modifica.
        \item Viene applicata la modifica.
    \end{enumerate}
    \item \textbf{Precondizioni}: Viene selezionata la funzione di modifica.
    \item \textbf{Postcondizioni}: Il testo presenta le modifiche.
    \item \textbf{Trigger}: L'utente si accorge di aver commesso un errore o vuole recuperare un'azione scartata.
\end{itemize}

% UC6.1 - Annullamento dell'ultima modifica (Undo)
\paragraph{UC 6.1: Annullamento dell'ultima modifica (Undo)}\mbox{}\

% use case diagram
\begin{figure}[H]
    \centering
    \includegraphics[width=0.6\textwidth]{gestione_modifiche_undo.png}
    \caption{UC6.1 – Annullamento dell'ultima modifica (Undo)}
\end{figure}

\begin{itemize}
    \item \textbf{Attore Principale}: Utente
    \item \textbf{Scenario Principale}: 
    \begin{enumerate}
        \item Esiste almeno una modifica annullabile.
        \item L'utente seleziona la funzione di undo.
        \item Viene annullata l'ultima modifica.
    \end{enumerate}
    \item \textbf{Precondizioni}: 
        \begin{itemize}
            \item Esiste almeno una modifica annullabile.
            \item Viene selezionata la funzione di undo.
        \end{itemize}
    \item \textbf{Postcondizioni}: L'ultima modifica viene annullata e il contenuto dell'editor viene aggiornato.
    \item \textbf{Trigger}: L'utente attiva il comando "Undo" per annullare l'ultima modifica apportata al testo.
\end{itemize}

% UC6.2 - Ripristino della modifica annullata (Redo)
\paragraph{UC 6.2: Ripristino della modifica annullata (Redo)}\mbox{}\

% use case diagram
\begin{figure}[H]
    \centering
    \includegraphics[width=0.6\textwidth]{gestione_modifiche_redo.png}
    \caption{UC6.2 – Ripristino della modifica annullata (Redo)}
\end{figure}

\begin{itemize}
    \item \textbf{Attore Principale}: Utente
    \item \textbf{Scenario Principale}: 
    \begin{enumerate}
        \item L'utente seleziona la funzione di ripristino (Redo).
        \item Il sistema riapplica l'ultima modifica che era stata annullata.
    \end{enumerate}
    \item \textbf{Precondizioni}: 
        \begin{itemize}
            \item Esiste almeno una modifica annullata precedentemente (tramite funzione Undo) che può essere ripristinata.
            \item L'utente attiva il comando di ripristino.
        \end{itemize}
    \item \textbf{Postcondizioni}: L'ultima modifica annullata viene ripristinata e il testo nell'editor viene aggiornato di conseguenza.
    \item \textbf{Trigger}: L'utente attiva il comando "Redo" per ripristinare una modifica precedentemente annullata.
\end{itemize}

% UC7 - Modifica modalità di visualizzazione
\subsubsection{UC 7: Modifica modalità di visualizzazione}

% use case diagram
\begin{figure}[H]
    \centering
    \includegraphics[width=0.6\textwidth]{modifica_visualizzazione.png}
    \caption{UC7 – Modifica modalità di visualizzazione}
\end{figure}

\begin{itemize}
    \item \textbf{Attore Principale}: Utente
    \item \textbf{Scenario Principale}: 
    \begin{enumerate}
        \item L'utente accede ai comandi per cambiare la modalità di visualizzazione dell'interfaccia.
        \item L'utente seleziona la modalità di visualizzazione desiderata.
        \item Il sistema riconfigura i pannelli dell'interfaccia in base alla scelta effettuata.
    \end{enumerate}
    \item \textbf{Precondizioni}: L'applicazione è avviata e l'editor è visualizzato a schermo.
    \item \textbf{Postcondizioni}: L'interfaccia utente si adatta mostrando i pannelli in base alla nuova modalità scelta.
    \item \textbf{Trigger}: L'utente vuole adattare l'interfaccia dell'editor alle proprie esigenze visive.
\end{itemize}

% UC7.1 - Visualizzazione solo Editor
\paragraph{UC 7.1: Visualizzazione solo Editor}\mbox{}\

% use case diagram
\begin{figure}[H]
    \centering
    \includegraphics[width=0.6\textwidth]{visualizzazione_solo_editor.png}
    \caption{UC7.1 – Visualizzazione solo Editor}
\end{figure}

\begin{itemize}
    \item \textbf{Attore Principale}: Utente
    \item \textbf{Scenario Principale}: 
    \begin{enumerate}
        \item L'utente seleziona la modalità di visualizzazione "Solo Editor".
    \end{enumerate}
    \item \textbf{Precondizioni}: L'utente accede ai comandi di modifica della visualizzazione.
    \item \textbf{Postcondizioni}: L'interfaccia mostra esclusivamente l'area di scrittura del codice Markdown.
    \item \textbf{Trigger}: L'utente vuole massimizzare lo spazio dedicato alla scrittura nascondendo l'anteprima.
\end{itemize}

% UC7.2 - Visualizzazione solo Render
\paragraph{UC 7.2: Visualizzazione solo Render}\mbox{}\

% use case diagram
\begin{figure}[H]
    \centering
    \includegraphics[width=0.6\textwidth]{visualizzazione_solo_render.png}
    \caption{UC7.2 – Visualizzazione solo Render}
\end{figure}

\begin{itemize}
    \item \textbf{Attore Principale}: Utente
    \item \textbf{Scenario Principale}: 
    \begin{enumerate}
        \item L'utente seleziona la modalità di visualizzazione "Solo Render".
    \end{enumerate}
    \item \textbf{Precondizioni}: L'utente accede ai comandi di modifica della visualizzazione.
    \item \textbf{Postcondizioni}: L'interfaccia mostra esclusivamente l'anteprima formattata (render) del documento.
    \item \textbf{Trigger}: L'utente vuole leggere il documento finale senza distrazioni dal codice Markdown.
\end{itemize}

% UC7.3 - Visualizzazione Split View
\paragraph{UC 7.3: Visualizzazione Split View}\mbox{}\

% use case diagram
\begin{figure}[H]
    \centering
    \includegraphics[width=0.6\textwidth]{visualizzazione_split_view.png}
    \caption{UC7.3 – Visualizzazione Split View}
\end{figure}

\begin{itemize}
    \item \textbf{Attore Principale}: Utente
    \item \textbf{Scenario Principale}: 
    \begin{enumerate}
        \item L'utente seleziona la modalità di visualizzazione "Split View".
    \end{enumerate}
    \item \textbf{Precondizioni}: L'utente accede ai comandi di modifica della visualizzazione.
    \item \textbf{Postcondizioni}: L'interfaccia si divide in due e mostra contemporaneamente e affiancati sia l'editor Markdown sia il render formattato.
    \item \textbf{Trigger}: L'utente vuole scrivere e contemporaneamente controllare in tempo reale l'anteprima grafica.
\end{itemize}


% --------------------------------------------------------------------------
% 2. Macroarea 2: Interazione AI (LLM)
% --------------------------------------------------------------------------
\newpage
\subsection{Macroarea 2 — Interazione AI (LLM)}
\textit{Elaborazione del testo tramite modelli linguistici e gestione degli output.}

% UC8 - Generazione riassunto del testo
\subsubsection{UC 8: Generazione riassunto del testo}

% use case diagram
\begin{figure}[H]
    \centering
    \includegraphics[width=0.6\textwidth]{riassunto_testo.png}
    \caption{UC8 – Riassunto del testo}
\end{figure}

\begin{itemize}
    \item \textbf{Attore Principale}: Utente
    \item \textbf{Attore Secondario}: LLM
    \item \textbf{Scenario Principale}: 
    \begin{enumerate}
        \item L'utente richiede un riassunto del testo selezionato.
        \item Il sistema invia la richiesta specifica all'LLM.
        \item L'LLM restituisce il riassunto richiesto.
    \end{enumerate}
    \item \textbf{Precondizioni}: L'applicazione è avviata, la connessione ai servizi AI è stabilita, una porzione di testo è selezionata nell'editor.
    \item \textbf{Postcondizioni}: Viene generato un riassunto che viene mostrato all'utente.
    \item \textbf{Trigger}: L'utente necessita di una sintesi rapida di un testo troppo lungo.
\end{itemize}

% UC9 - Traduzione del testo
\subsubsection{UC 9: Traduzione del testo}

% use case diagram
\begin{figure}[H]
    \centering
    \includegraphics[width=0.6\textwidth]{traduzione_testo.png}
    \caption{UC9 – Traduzione del testo}
\end{figure}

\begin{itemize}
    \item \textbf{Attore Principale}: Utente
    \item \textbf{Attore Secondario}: LLM
    \item \textbf{Scenario Principale}: 
    \begin{enumerate}
        \item L'utente richiede la traduzione del testo selezionato nella lingua scelta.
        \item Il sistema invia la richiesta specifica all'LLM.
        \item L'LLM restituisce il testo tradotto richiesto.
    \end{enumerate}
    \item \textbf{Precondizioni}: L'applicazione è avviata, la connessione ai servizi AI è stabilita, una porzione di testo è selezionata nell'editor.
    \item \textbf{Postcondizioni}: Viene generato una traduzione che viene mostrato all'utente.
    \item \textbf{Trigger}: L'utente ha bisogno di convertire una porzione di testo in una lingua differente.
\end{itemize}
% UC10 - Riscrittura del testo
\subsubsection{UC 10: Riscrittura del testo}

% use case diagram
\begin{figure}[H]
    \centering
    \includegraphics[width=0.6\textwidth]{riscrittura_testo.png}
    \caption{UC10 – Riscrittura del testo}
\end{figure}

\begin{itemize}
    \item \textbf{Attore Principale}: Utente
    \item \textbf{Attore Secondario}: LLM
    \item \textbf{Scenario Principale}: 
    \begin{enumerate}
        \item L'utente richiede una riscrittura del testo selezionato.
        \item Il sistema invia la richiesta specifica all'LLM.
        \item L'LLM restituisce la riscrittura richiesta.
    \end{enumerate}
    \item \textbf{Precondizioni}: L'applicazione è avviata, la connessione ai servizi AI è stabilita, una porzione di testo è selezionata nell'editor.
    \item \textbf{Postcondizioni}: Viene generato una rielaborazione del testo che viene mostrata all'utente.
    \item \textbf{Trigger}: L'utente non è soddisfatto della forma del testo e vuole che l'AI lo rielabori.
\end{itemize}
% UC11 - Generazione automatica da prompt (Distant Writing)
\subsubsection{UC 11: Generazione automatica da prompt (Distant Writing)}

% use case diagram
\begin{figure}[H]
    \centering
    \includegraphics[width=0.6\textwidth]{11_generazione_automatica_da_prompt_distant_writing.png}
    \caption{UC11 – Generazione automatica da prompt (Distant Writing)}
\end{figure}

\begin{itemize}
    \item \textbf{Attore Principale}: Utente
    \item \textbf{Attore Secondario}: LLM
    \item \textbf{Scenario Principale}: 
    \begin{enumerate}
        \item L'utente inserisce un'istruzione testuale (prompt) in un apposito campo.
        \item Il sistema invia il prompt all'LLM.
        \item L'LLM genera il testo e lo restituisce al sistema.
        \item Il sistema fornisce il testo generato all'utente (gestito tramite UC 12).
    \end{enumerate}
    \item \textbf{Precondizioni}: L'applicazione è avviata e la connessione ai servizi AI è stabilita.
    \item \textbf{Postcondizioni}: Viene generata una proposta di testo basata sul prompt fornito.
    \item \textbf{Trigger}: L'utente vuole che l'AI generi nuovi contenuti testuali a partire da una sua istruzione diretta.
\end{itemize}

% UC12 - Critica del testo (Sei Cappelli per Pensare)
\subsubsection{UC 12: Critica del testo (Sei Cappelli per Pensare)}

% use case diagram
\begin{figure}[H]
    \centering
    \includegraphics[width=0.6\textwidth]{12_critica_del_testo_sei_cappelli_per_pensare.png}
    \caption{UC12 – Critica del testo (Sei Cappelli per Pensare)}
\end{figure}

\begin{itemize}
    \item \textbf{Attore Principale}: Utente
    \item \textbf{Attore Secondario}: LLM
    \item \textbf{Scenario Principale}: 
    \begin{enumerate}
        \item L'utente richiede un'analisi critica del testo selezionato.
        \item L'utente sceglie una delle sei modalità (Cappelli).
        \item Il sistema invia la richiesta specifica all'LLM.
        \item L'LLM restituisce il feedback critico.
    \end{enumerate}
    \item \textbf{Precondizioni}: L'applicazione è avviata, la connessione ai servizi AI è stabilita, una porzione di testo è selezionata nell'editor.
    \item \textbf{Postcondizioni}: Viene generata una proposta di analisi che viene mostrata all'utente.
    \item \textbf{Trigger}: L'utente desidera che il testo venga analizzato criticamente adottando una specifica prospettiva.
\end{itemize}

% UC12.1 - Analisi Cappello Bianco (Dati e fatti)
\paragraph{UC 12.1: Analisi Cappello Bianco (Dati e fatti)}\mbox{}\

% use case diagram
\begin{figure}[H]
    \centering
    \includegraphics[width=0.6\textwidth]{12-1_analisi_cappello_bianco_dati_e_fatti.png}
    \caption{UC12.1 – Analisi Cappello Bianco (Dati e fatti)}
\end{figure}

\begin{itemize}
    \item \textbf{Attore Principale}: Utente
    \item \textbf{Attore Secondario}: LLM
    \item \textbf{Scenario Principale}: 
    \begin{enumerate}
        \item L'utente richiede un'analisi oggettiva.
        \item Il sistema istruisce l'LLM a evidenziare solo informazioni certe, dati e lacune informative nel testo selezionato.
        \item L'LLM restituisce il feedback critico.
    \end{enumerate}
    \item \textbf{Precondizioni}: L'applicazione è avviata, la connessione ai servizi AI è stabilita, una porzione di testo è selezionata nell'editor.
    \item \textbf{Postcondizioni}: Viene visualizzata un'analisi critica basata su fatti e dati oggettivi.
    \item \textbf{Trigger}: L'utente seleziona il Cappello Bianco per verificare i fatti, i dati e l'oggettività del testo.
\end{itemize}

% UC12.2 - Analisi Cappello Rosso (Emozioni)
\paragraph{UC 12.2: Analisi Cappello Rosso (Emozioni)}\mbox{}\

% use case diagram
\begin{figure}[H]
    \centering
    \includegraphics[width=0.6\textwidth]{12-2_analisi_cappello_rosso_emozioni.png}
    \caption{UC12.2 – Analisi Cappello Rosso (Emozioni)}
\end{figure}

\begin{itemize}
    \item \textbf{Attore Principale}: Utente
    \item \textbf{Attore Secondario}: LLM
    \item \textbf{Scenario Principale}: 
    \begin{enumerate}
        \item L'utente richiede un'analisi emotiva.
        \item L'LLM analizza i sentimenti, le intuizioni e le reazioni viscerali trasmesse dal testo.
        \item L'LLM restituisce il feedback critico.
    \end{enumerate}
    \item \textbf{Precondizioni}: L'applicazione è avviata, la connessione ai servizi AI è stabilita, una porzione di testo è selezionata nell'editor.
    \item \textbf{Postcondizioni}: Viene visualizzata un'analisi del tono emotivo del testo.
    \item \textbf{Trigger}: L'utente seleziona il Cappello Rosso per valutare l'impatto emotivo e viscerale del testo.
\end{itemize}

% UC12.3 - Analisi Cappello Nero (Rischi)
\paragraph{UC 12.3: Analisi Cappello Nero (Rischi)}\mbox{}\

% use case diagram
\begin{figure}[H]
    \centering
    \includegraphics[width=0.6\textwidth]{12-3_analisi_cappello_nero_rischi.png}
    \caption{UC12.3 – Analisi Cappello Nero (Rischi)}
\end{figure}

\begin{itemize}
    \item \textbf{Attore Principale}: Utente
    \item \textbf{Attore Secondario}: LLM
    \item \textbf{Scenario Principale}: 
    \begin{enumerate}
        \item L'utente richiede una valutazione critica negativa.
        \item L'LLM individua potenziali punti deboli, rischi o incongruenze logiche.
        \item L'LLM restituisce il feedback critico.
    \end{enumerate}
    \item \textbf{Precondizioni}: L'applicazione è avviata, la connessione ai servizi AI è stabilita, una porzione di testo è selezionata nell'editor.
    \item \textbf{Postcondizioni}: Viene visualizzata un'analisi delle criticità del testo.
    \item \textbf{Trigger}: L'utente seleziona il Cappello Nero per individuare criticità, rischi e punti deboli.
\end{itemize}

% UC12.4 - Analisi Cappello Giallo (Opportunità)
\paragraph{UC 12.4: Analisi Cappello Giallo (Opportunità)}\mbox{}\

% use case diagram
\begin{figure}[H]
    \centering
    \includegraphics[width=0.6\textwidth]{12-4_analisi_cappello_giallo_opportunita.png}
    \caption{UC12.4 – Analisi Cappello Giallo (Opportunità)}
\end{figure}

\begin{itemize}
    \item \textbf{Attore Principale}: Utente
    \item \textbf{Attore Secondario}: LLM
    \item \textbf{Scenario Principale}: 
    \begin{enumerate}
        \item L'utente richiede una valutazione ottimistica.
        \item L'LLM mette in luce i benefici, i punti di forza e il valore aggiunto del testo.
        \item L'LLM restituisce il feedback critico.
    \end{enumerate}
    \item \textbf{Precondizioni}: L'applicazione è avviata, la connessione ai servizi AI è stabilita, una porzione di testo è selezionata nell'editor.
    \item \textbf{Postcondizioni}: Viene visualizzata un'analisi dei vantaggi e delle opportunità.
    \item \textbf{Trigger}: L'utente seleziona il Cappello Giallo per scoprire i vantaggi, i benefici e le opportunità espresse.
\end{itemize}

% UC12.5 - Analisi Cappello Verde (Creatività)
\paragraph{UC 12.5: Analisi Cappello Verde (Creatività)}\mbox{}\

% use case diagram
\begin{figure}[H]
    \centering
    \includegraphics[width=0.6\textwidth]{12-5_analisi_cappello_verde_creativita.png}
    \caption{UC12.5 – Analisi Cappello Verde (Creatività)}
\end{figure}

\begin{itemize}
    \item \textbf{Attore Principale}: Utente
    \item \textbf{Attore Secondario}: LLM
    \item \textbf{Scenario Principale}: 
    \begin{enumerate}
        \item L'utente richiede alternative creative.
        \item L'LLM propone nuove idee, approcci originali o variazioni fantasiose sul tema trattato.
        \item L'LLM restituisce il feedback critico.
    \end{enumerate}
    \item \textbf{Precondizioni}: L'applicazione è avviata, la connessione ai servizi AI è stabilita, una porzione di testo è selezionata nell'editor.
    \item \textbf{Postcondizioni}: Vengono visualizzati suggerimenti creativi e alternative al testo.
    \item \textbf{Trigger}: L'utente seleziona il Cappello Verde per cercare idee alternative, spunti creativi e innovazioni.
\end{itemize}

% UC12.6 - Analisi Cappello Blu (Logica/Controllo)
\paragraph{UC 12.6: Analisi Cappello Blu (Logica/Controllo)}\mbox{}\

% use case diagram
\begin{figure}[H]
    \centering
    \includegraphics[width=0.6\textwidth]{12-6_analisi_cappello_blu_logica_controllo.png}
    \caption{UC12.6 – Analisi Cappello Blu (Logica/Controllo)}
\end{figure}

\begin{itemize}
    \item \textbf{Attore Principale}: Utente
    \item \textbf{Attore Secondario}: LLM
    \item \textbf{Scenario Principale}: 
    \begin{enumerate}
        \item L'utente richiede un'analisi strutturale.
        \item L'LLM valuta l'organizzazione del pensiero, la coerenza complessiva e il processo logico seguito.
        \item L'LLM restituisce il feedback critico.
    \end{enumerate}
    \item \textbf{Precondizioni}: L'applicazione è avviata, la connessione ai servizi AI è stabilita, una porzione di testo è selezionata nell'editor.
    \item \textbf{Postcondizioni}: Viene visualizzata una sintesi del processo logico e della struttura del testo.
    \item \textbf{Trigger}: L'utente seleziona il Cappello Blu per analizzare la struttura logica e l'organizzazione del documento.
\end{itemize}

% UC13 - Gestione proposta AI
\subsubsection{UC 13: Gestione proposta AI}

% use case diagram
\begin{figure}[H]
    \centering
    \includegraphics[width=0.6\textwidth]{13_gestione_proposta_ai.png}
    \caption{UC13 – Gestione proposta AI}
\end{figure}

\begin{itemize}
    \item \textbf{Attore Principale}: Utente
    \item \textbf{Scenario Principale}: 
    \begin{enumerate}
        \item L'utente valuta la proposta fornita dal sistema.
        \item L'utente decide se accettare (UC 12.1) o rifiutare (UC 12.2) la proposta.
    \end{enumerate}
    \item \textbf{Precondizioni}: L'LLM ha restituito una proposta (testo generato, traduzione o analisi) che è visualizzata a schermo.
    \item \textbf{Postcondizioni}: La visualizzazione della proposta viene chiusa e l'editor viene aggiornato o resta invariato.
    \item \textbf{Trigger}: Il sistema ha terminato l'elaborazione tramite LLM e mostra il risultato all'utente.
\end{itemize}

% UC13.1 - Inserimento proposta nel documento
\paragraph{UC 13.1: Inserimento proposta nel documento}\mbox{}\

% use case diagram
\begin{figure}[H]
    \centering
    \includegraphics[width=0.6\textwidth]{13-1_inserimento_proposta_nel_documento.png}
    \caption{UC13.1 – Inserimento proposta nel documento}
\end{figure}

\begin{itemize}
    \item \textbf{Attore Principale}: Utente
    \item \textbf{Scenario Principale}: 
    \begin{enumerate}
        \item L'utente accetta tramite bottone l'inserimento della proposta fornita.
        \item Il sistema inserisce nel testo l'output dell'AI.
    \end{enumerate}
    \item \textbf{Precondizioni}: Una proposta AI è visibile a schermo.
    \item \textbf{Postcondizioni}: Il contenuto dell'editor è aggiornato con il nuovo testo.
    \item \textbf{Trigger}: L'utente ritiene valida la proposta dell'AI e preme il tasto di conferma.
\end{itemize}

% UC13.2 - Rifiuto della proposta AI
\paragraph{UC 13.2: Rifiuto della proposta AI}\mbox{}\

% use case diagram
\begin{figure}[H]
    \centering
    \includegraphics[width=0.6\textwidth]{13-2_rifiuto_della_proposta_ai.png}
    \caption{UC13.2 – Rifiuto della proposta AI}
\end{figure}

\begin{itemize}
    \item \textbf{Attore Principale}: Utente
    \item \textbf{Scenario Principale}: 
    \begin{enumerate}
        \item L'utente rifiuta tramite bottone l'inserimento della proposta fornita.
        \item Il sistema scarta l'output dell'AI senza modificare il documento originale.
    \end{enumerate}
    \item \textbf{Precondizioni}: Una proposta AI è visibile a schermo.
    \item \textbf{Postcondizioni}: La proposta viene eliminata e il testo nell'editor rimane inalterato.
    \item \textbf{Trigger}: L'utente ritiene la proposta inadeguata e preme il tasto per annullarla e chiuderla.
\end{itemize}

% UC14 - Gestione errori di sistema LLM
\subsubsection{UC 14: Gestione errori di sistema LLM}

% use case diagram
\begin{figure}[H]
    \centering
    \includegraphics[width=0.6\textwidth]{14_gestione_errori_di_sistema_llm.png}
    \caption{UC14 – Gestione errori di sistema LLM}
\end{figure}

\begin{itemize}
    \item \textbf{Attore Principale}: Utente
    \item \textbf{Attore Secondario}: LLM
    \item \textbf{Scenario Principale}: 
    \begin{enumerate}
        \item Il sistema invia una richiesta all'LLM.
        \item L'LLM restituisce un errore (es: timeout, superamento limiti API, assenza di rete).
        \item Il sistema intercetta l'errore e mostra un messaggio informativo all'utente.
    \end{enumerate}
    \item \textbf{Precondizioni}: È stata avviata un'interazione che richiede l'uso dell'AI.
    \item \textbf{Postcondizioni}: L'utente è informato del fallimento e il sistema ripristina lo stato precedente all'operazione.
    \item \textbf{Trigger}: La comunicazione con il servizio AI fallisce, scade per timeout o restituisce un errore.
\end{itemize}

% UC15 - Interazione con ChatBot
\subsubsection{UC 15: (Opzionale) Interazione con ChatBot}

% use case diagram
\begin{figure}[H]
    \centering
    \includegraphics[width=0.6\textwidth]{Interazione_con_ChatBot.png}
    \caption{UC15 – Interazione con ChatBot}
\end{figure}


\begin{itemize}
    \item \textbf{Attore Principale}: Utente
    \item \textbf{Attore Secondario}: LLM
    \item \textbf{Scenario Principale}: 
    \begin{enumerate}
        \item L'utente apre l'interfaccia del ChatBot.
        \item L'utente invia un messaggio di testo o seleziona un comando predefinito.
        \item Il sistema inoltra la richiesta al modello linguistico (LLM).
        \item Il sistema visualizza la risposta generata dall'AI nella chat.
    \end{enumerate}
    \item \textbf{Precondizioni}: Il sistema è connesso a internet e il servizio LLM è disponibile.
    \item \textbf{Postcondizioni}: L'utente riceve una risposta coerente con il messaggio inviato.
    \item \textbf{Trigger}: L'utente ha bisogno di dialogare in maniera conversazionale con l'assistente virtuale.
\end{itemize}


% --------------------------------------------------------------------------
% 3. Macroarea 3: Gestione I/O e Persistenza
% --------------------------------------------------------------------------
\newpage
\subsection{Macroarea 3 — Gestione I/O e Persistenza}
\textit{Salvataggio, caricamento ed esportazione.}

% UC16 - Download della nota in formato Markdown
\subsubsection{UC 16: Download della nota in formato Markdown}

% use case diagram
\begin{figure}[H]
    \centering
    \includegraphics[width=0.6\textwidth]{Download_nota_formato_Markdown.png}
    \caption{UC16 – Download della nota in formato Markdown}
\end{figure}

\begin{itemize}
    \item \textbf{Attore Principale}: Utente
    \item \textbf{Scenario Principale}: 
    \begin{enumerate}
        \item L'utente seleziona la funzione di download/salvataggio locale.
        \item Il sistema genera un file con estensione .md contenente il testo dell'editor.
        \item Il sistema permette all'utente di scegliere la destinazione sul proprio dispositivo.
    \end{enumerate}
    \item \textbf{Precondizioni}: L'editor contiene del testo (anche vuoto).
    \item \textbf{Postcondizioni}: Il file Markdown viene salvato nel file system locale dell'utente.
    \item \textbf{Trigger}: L'utente vuole salvare una copia locale del file sorgente .md.
\end{itemize}

% UC17 - Upload di una nota locale
\subsubsection{UC 17: Upload di una nota locale}

% use case diagram
\begin{figure}[H]
    \centering
    \includegraphics[width=0.6\textwidth]{Upload_nota_locale.png}
    \caption{UC17 – Upload di una nota locale}
\end{figure}

\begin{itemize}
    \item \textbf{Attore Principale}: Utente
    \item \textbf{Scenario Principale}: 
    \begin{enumerate}
        \item L'utente attiva la funzione di caricamento file.
        \item L'utente seleziona un file .md dal proprio dispositivo.
        \item Il sistema legge il contenuto del file e lo carica nell'editor, sovrascrivendo o aggiungendo al testo presente.
    \end{enumerate}
    \item \textbf{Precondizioni}: L'utente possiede un file Markdown compatibile.
    \item \textbf{Postcondizioni}: Il contenuto del file selezionato è visibile e modificabile nell'editor.
    \item \textbf{Trigger}: L'utente vuole caricare nel sistema un file Markdown già esistente nel proprio computer.
\end{itemize}


% UC18 - Esportazione della nota (PDF, HTML, JSON)
\subsubsection{UC 18: Esportazione della nota (PDF, HTML, JSON)}

% use case diagram
\begin{figure}[H]
    \centering
    \includegraphics[width=0.6\textwidth]{esportazione_nota.png}
    \caption{UC18 – Esportazione della nota}
\end{figure}

\begin{itemize}
    \item \textbf{Attore Principale}: Utente
    \item \textbf{Scenario Principale}: 
    \begin{enumerate}
        \item L'utente seleziona la funzione di esportazione.
        \item L'utente sceglie il formato desiderato (PDF, HTML o JSON).
        \item Il sistema converte il contenuto Markdown nel formato scelto e avvia il download.
    \end{enumerate}
    \item \textbf{Precondizioni}: Il sistema è in grado di eseguire la conversione del formato.
    \item \textbf{Postcondizioni}: L'utente ottiene una copia del documento nel formato non-Markdown selezionato.
    \item \textbf{Trigger}: L'utente desidera esportare il documento compilato in un formato idoneo alla distribuzione.
\end{itemize}

% UC19 - Gestione Persistenza su Database remoto
\subsubsection{UC 19: (Opzionale) Gestione Persistenza su Database remoto}

% use case diagram riassuntivo per la persistenza
\begin{figure}[H]
    \centering
    \includegraphics[width=0.7\textwidth]{gestione_persistenza.png}
    \caption{UC19 – Gestione Persistenza su Database remoto}
\end{figure}

\begin{itemize}
    \item \textbf{Attore Principale}: Utente
    \item \textbf{Attore Secondario}: Database (Sistema esterno)
    \item \textbf{Descrizione}: Il sistema permette all'utente autenticato di gestire le proprie note su uno spazio di archiviazione remoto, garantendo la disponibilità dei dati su diversi dispositivi.
    \item \textbf{Precondizioni}: L'utente deve aver effettuato l'accesso (Login) al sistema.
    \item \textbf{Postcondizioni}: I cambiamenti effettuati (salvataggio, caricamento o collegamento) sono sincronizzati con il database.
    \item \textbf{Trigger}: L'utente vuole gestire i propri file sfruttando l'archiviazione cloud del sistema.
\end{itemize}

% UC19.1 - Salvataggio nota su database
\paragraph{UC 19.1: Salvataggio nota su database}
\begin{itemize}
    \item \textbf{Scenario Principale}: L'utente attiva il comando di salvataggio remoto; il sistema invia il contenuto dell'editor al database associandolo al profilo dell'utente.
    \item \textbf{Trigger}: L'utente attiva il comando per sincronizzare la nota corrente sul server remoto.
\end{itemize}

% UC19.2 - Caricamento nota da database
\paragraph{UC 19.2: Caricamento nota da database}
\begin{itemize}
    \item \textbf{Scenario Principale}: L'utente richiede la lista delle note salvate; il sistema recupera i metadati dal database e, dopo la selezione, scarica il testo nell'editor.
    \item \textbf{Trigger}: L'utente decide di aprire una nota precedentemente salvata nel proprio profilo cloud.
\end{itemize}

% UC19.3 - Creazione collegamento tra note nel database
\paragraph{UC 19.3: Creazione collegamento tra note nel database}
\begin{itemize}
    \item \textbf{Scenario Principale}: L'utente inserisce un riferimento (ID o Link) a una nota già presente nel database per creare una connessione logica tra i documenti.
    \item \textbf{Trigger}: L'utente vuole collegare la nota corrente a un'altra nota nel database creando un link bidirezionale.
\end{itemize}